% Part: history
% Chapter: biographies 
% Section: alonzo-church

\documentclass[../../../include/open-logic-section]{subfiles}

\begin{document}

\olfileid{his}{bio}{chu}

\olsection{ॲलोन्झो चर्च}

\olphoto{church-alonzo}{ॲलोन्झो चर्च}

ॲलोन्झो चर्च यांचा जन्म १४ जून १९०३ रोजी वॉशिंग्टन, डीसी येथे झाला. लहानपणी हवेच्या बंदुकीशी संबंधित एका अपघातामुळे त्यांचा एक डोळा आंधळा झाला. कनेक्टिकट येथील महाविद्यालयपूर्व शाळेचे शिक्षण त्यांनी १९२० मध्ये पूर्ण केले आणि त्याच वर्षी प्रिन्स्टन विद्यापीठात शिक्षण सुरू केले. १९२७ मध्ये त्यांनी डॉक्टरेट पूर्ण केली. दोन वर्षे परदेशात राहिल्यानंतर चर्च प्रिन्स्टनला परतले. ते अत्यंत नम्र आणि काळजीपूर्वक काम करणारे म्हणून ओळखले जात. फळ्यावरचे त्यांचे लेखन निर्दोष असे. महत्त्वाचे कागद जतन करण्यासाठी ते त्यांवर ड्यूको सिमेंट या पारदर्शक गोंदाचा पातळ थर लावत. शैक्षणिक कामाव्यतिरिक्त त्यांना विज्ञानकथा मासिके वाचायला आवडत; कथांमध्ये काही चूक आढळली तर संपादकांना पत्र लिहायलाही ते कचरत नसत.

चर्च यांचे शैक्षणिक योगदान मोठे होते. त्यांचे विद्यार्थी स्टीफन क्लिनी आणि बार्कली रॉसर यांच्यासह त्यांनी प्रभावी गणनीयतेची एक उपपत्ती, म्हणजे लॅम्डा कलन, विकसित केले. हे काम ॲलन ट्यूरिंग यांनी ट्यूरिंग यंत्र विकसित केले त्या कामापासून स्वतंत्रपणे झाले. गणनीयतेच्या या दोन व्याख्या सममूल्य आहेत. त्यांतून आज \emph{चर्च–ट्यूरिंग प्रबंध} म्हणून ओळखले जाणारे प्रतिपादन पुढे येते: नैसर्गिक संख्यांवरील एखादे फलन प्रभावीपणे गणनीय असते जर आणि फक्त जर ते ट्यूरिंग यंत्राने (किंवा लॅम्डा कलनाने) गणनीय असेल. चर्च यांनी आज \emph{चर्च यांचे प्रमेय} म्हणून ओळखला जाणारा निष्कर्षही सिद्ध केला: प्रथम-क्रम सूत्रांच्या वैधतेची निर्णय समस्या सोडवता येत नाही.

चर्च यांनी वृद्धापकाळापर्यंत काम चालू ठेवले. १९६७ मध्ये त्यांनी प्रिन्स्टन सोडून यूसीएलए येथे प्रवेश केला. १९९० मध्ये निवृत्त होईपर्यंत ते तेथे प्राध्यापक होते. ११ ऑगस्ट १९९५ रोजी वयाच्या ९२व्या वर्षी त्यांचे निधन झाले.

\begin{reading}
चर्च यांचे संक्षिप्त चरित्र \citet{EndertonND} येथे पाहा. लॅम्डा कलन, चर्च यांचा प्रबंध आणि निर्णय समस्या यांवरील चर्च यांचे मूळ लेख \citet{Church1936,Church1936a} येथे आहेत. १९३०च्या दशकातील प्रिन्स्टनच्या गणिती समुदायाविषयी चर्च यांची मुलाखत \citet{Aspray1984} यांनी नोंदवली आहे. चर्च यांनी १९३६ ते १९७९ या काळात \emph{Journal of Symbolic Logic} साठी अनेक पुस्तक परीक्षणे लिहिली. ती सर्व जॉन मॅकफार्लेन यांच्या संकेतस्थळावर संग्रहित आहेत \citep{MacFarlane2015}.
\end{reading} 
\end{document}
